\subsection{The m-strain block GLV model with  rank one invading blocks}
One of the most convincing results for the utility of CRN-ME tools is the
Lyapunov Ansatz below for a  model with  $m$ competing irreducible
rank-one blocks, and with no direct block-to-block coupling in the
transversal equations:
\be{eq:rankone-m}
\begin{cases}
s'
=
\Lambda-\mu s-s\displaystyle\sum_{k=1}^m\phi_k,
\\[1mm]
z_k'
=
\left(sw_k\pi_k^\top-V_k\right)z_k,
\qquad
k=1,\ldots,m,
\\[1mm]
\phi_k
=
\pi_k^\top z_k.
\end{cases}
\ee
Here $z_k\in\R_{\ge0}^{n_k}$, $w_k\gg0$, $\pi_k\gg0$,
${\bf 1}^\top\pi_k=1$, and    $V_k$ are  nonsingular
$M$-matrices, with  $V_k^{-1}\gg0$ (internal transitions/removal).

Note the ``structure aware" writing of the ODE, which partitions the variables into  $m$ non-interacting disjoint minimal siphon/maximal invariant faces/``vector Lotka-Volterra"/invading  groups, and a non-siphon/non-DFE variable $s$, and the rank-one regular splitting of the transversal Jacobian/invading matrices for each block:
\[
J_k^\perp(s)=F_k(s)-V_k,
\qquad
F_k(s)=sw_k\pi_k^\top\ \gg0(\text{rank one, new infections}).
\]


We assume that
\be{wr}
{\bf 1}^\top V_k\geq\mu{\bf 1}^\top,
\qquad
k=1,\ldots,m
\ee
(in epidemiologic terms, the decay of the invading strain compartments surpasses
that of the input, ensuring the invariance of the set below).
Then
\[
\Omega
=
\left\{
s\geq0,\ z_k\geq0:
s+\sum_{k=1}^m{\bf 1}^\top z_k\leq s_0
=\frac{\Lambda}{\mu}
\right\}
\]
is forward invariant.

The monograph considers  variations of this model. When all blocks have dimension 1, we have a subclass of GLV models, studied in Part I (without the non-interaction assumption).  Part II considers the cases when $V_k$ are diagonal, and also the general case of  
of Metzler-Hurwicz   $-V_k$, but only when $m=2$. Finally, Part III considers a model which extends the models of Recker-Gupta and Smith-Browne in immuno-virology.

